\SourceLeaf{480}

\Needspace{5\baselineskip}\begin{center}ANHANG.\end{center}
\BookRule{0.15\linewidth}

\Needspace{5\baselineskip}\begin{center}\small HANDSCHRIFTLICHE AUFZEICHNUNGEN VON GAUSS.\end{center}


Zu Art. 40. \textit{Accedente numero tertio \(C\), sit \(\lambda'\) maximus divisor communis numerorum \(\lambda,C\), determinenturque numeri \(k,\gamma\) ita ut sit \(k\lambda+\gamma C=\lambda'\), unde erit \(k\alpha A+k\beta B+\gamma C=\lambda'\). Manifesto autem \(\lambda'\) est divisor communis numerorum \(A,B.C\), et quidem maximus, si enim exstaret major \(=\theta\), foret \[k\alpha.\frac{A}{\theta}+k\beta.\frac{B}{\theta}+\gamma.\frac{C}{\theta}=\frac{\lambda'}{\theta}\quad\textit{integer, Q. E. A.}\]}


\textit{Factum est itaque quod propositum fuerat. dum statuimus \(k\alpha=a,\ k\beta=b,\ \gamma=c,\ \lambda'=\mu\).}


Zu Art. 114. \textit{Concinnius.demonstratio ita adornatur \((a^{3n}-a^n)^2=2+(a^{4n}+1)(a^{2n}-2)\), \((a^{3n}+a^n)^2=-2+(a^{4n}+1)(a^{2n}+2)\) adeoque \(\surd2\equiv\pm(a^{3n}-a^n),\ \surd-2\equiv\pm(a^{3n}+a^n)\ (\mathrm{mod.}\ 8n+1)\)}


Zu den in Art. 256 VI angegebenen 16 positiven Determinanten von der Form \(8n+5\), für welche die Anzahl der eigentlich primitiven Classen drei mal grösser ist als die der uneigentlich primitiven ‘37, .. 573’ sind hinzugefügt: \textit{677, 701, 709, 757, 781, 813, 829, 877, 885, 901, 909, 925, 933, 973, 997 adeoque inter 125 exstant 31.}


Zu den Worten des Art. 301: Hoc modo summa mult. gen. pro dett. \(-\): usque ad \(-100\) invenitur \(=234,4\) quum revera sit 233. — — \textit{a \(-1\) usque ad \(-3000\), tabula 11166, formula 11167,9.}


Zu Art. 336. \textit{Wären alle Zahlen der Form \(2^{2^m}+1\) Primzahlen, so würde ein hinlänglich genäherter Ausdruck für die Menge der in Rede stehenden Zahlen \((N)\) kleiner als die gegebene Zahl \(M\), folgender sein \(\frac12\left(\frac{\log M}{\log2}\right)^2\).}


Zu dem Lehrsatze in Art. 42 über die Theiler einer algebraischen ganzen rationalen Function mit ganzzahligen Coëfficienten — — \textit{1797 Jul.22.}


Zu den Worten des Art. 130: Postquam rigorose demonstravimus, quemvis numerum primum formae \(4n+1\), et positive et negative acceptum, alicuius numeri primi ipso minoris non-residuum esse, — — \textit{Hanc demonstrationem deteximus 1796 Apr. 8.}


Zu Art. 131. \textit{Theorema fundamentale per inductionem detectum 1795 Martio. Demonstratio prima, quae in hac sectione traditur, inventa 1796 Apr.}

\SourceLeaf{481}


Zu den Worten des Art. 133: Investigationem (theor. fund.) adhuc generalius instituamus. Contemplemur duos numeros quoscunque impares inter se primos, signis quibuscunque affectos, \(P\) et \(Q\). — — \textit{1796 Apr. 29.}


Zu den Worten des Art. 145: Praeterea theoremata ad residua \(+2\) et \(-2\) pertinentia tunc supponi debuissent; quum vero nostra demonstratio absque his theorematibus sit perfecta, novam hinc methodum nanciscimur, illa demonstrandi. — — \textit{1797 Febr. 4.}


Zu der Ueberschrift der Sectio quinta: De formis aequationibusque indeterminatis secundi gradus. — — \textit{Inde a Jun. 22. 1796.}


Zu den Worten des Art. 234: .. ad aliud argumentum gravissimum transimus a nemine hucusque attactum, de formarum compositione. — — \textit{Hae disquiss. inchoatae autumno 1798.}


Zu den Worten des Art. 262: Ex hoc principio methodum novam haurire possumus, non modo theorema fundamentale, sed etiam reliqua theoremata Sect. praec. ad residua \(-1,\allowbreak +2,\allowbreak -2\) pertinentia demonstrandi, — — \textit{Principia huius methodi primum se obtulerunt 1796 Jul. 27, at exculta et ad formam praesentem redacta Vere a. 1800.}


Zu den Worten des Art. 266: .. sed quoniam complures veritates ad has spectantes, eaeque pulcherrimae, adhuc supersunt, quarum fons proprius in theoria formarum ternariarum secundi gradus est quaerendus, brevem ad hanc theoriam digressionem hic intercalamus, — — \textit{Febr. 14. 1799.}


Zu den Worten des Art. 272: scilicet ostendendo, primo, quo pacto quaevis forma ternaria ad formam simpliciorem reduci possit, dein, formarum simplicissimarum (ad quas per tales reductiones perveniatur), multitudinem pro quovis determinante dato esse finitam. — — \textit{1800 Febr. 13.}


Zu den Worten des Art. 287 III: Prorsus simili ratione probatur, in ordine improprie primitivo eos characteres, qui per praecepta art. 264 II, III soli possibiles inveniuntur, omnes possibiles esse, sive sint \(P\) sive \(Q\). — Haecce theoremata, etc. — — \textit{Demonstratione primum munita sunt Mense Aprili 1798.}


Zu den Worten des Art. 302: Multitudo classium mediocris autem (quae definitione opus non habebit) valde regulariter crescit, — — \textit{Idea prima initio a. 1799.}


Zu den Worten des Art. 306 X: Denique observamus, quum omnes proprietates in hoc art. et praec. consideratae imprimis a numero \(n\) pendeant, qui simile quid est ac \(p-1\) in Sect. III, hunc numerum summa attentione dignum esse; quamobrem quam maxime optandum esset, ut inter ipsum atque determinantem, ad quem pertinet, nexus generalis detegatur. — — \textit{Ex voto nobis-ac sic successit ut nihil amplius desiderandum supersit Nov. 30—Dec. 3. 1800.}


Zu Art. 365. \textit{Circulum in 17 partes divisibilem esse geometrice, deteximus 1796 Mart. 30.}

\BookRule{0.3\linewidth}

\clearpage
